\section{Continuous spectrum}

\SourcePage{29}{32}
All the relations obtained in \secref{3}, \secref{4} for the properties of discrete-spectrum
eigenfunctions extend without difficulty to a continuous spectrum of
eigenvalues.

Let $f$ be a physical quantity possessing a continuous spectrum. We shall
denote its eigenvalues simply by the same letter $f$ without an index; the
corresponding eigenfunctions will be written $\Psi_f$. In the same way that an
arbitrary wave function $\Psi$ admits the series expansion (3.2) in
eigenfunctions of a discrete-spectrum quantity, it can also be expanded --- this
time as an integral --- over the complete system of eigenfunctions of a
continuous-spectrum quantity. Such an expansion takes the form
\begin{equation}
  \Psi(q)=\int a_f\Psi_f(q)\,df,
  \tag{5.1}
\end{equation}
where the integration is taken over the whole domain of values that the
quantity $f$ can assume.

More complicated than in the discrete-spectrum case is the question of
normalizing the continuous-spectrum eigenfunctions. As we shall see below, the
requirement that the integral of the squared modulus equal unity cannot be
satisfied here. Instead, let us set ourselves
the aim of normalizing the functions $\Psi_f$ in such a way that $|a_f|^2\,df$
represents the probability for the physical quantity under consideration, in
the state described by the wave function $\Psi$, to have a value in the given
interval between $f$ and $f+df$.
\SourcePage{30}{33}
Since the sum of the probabilities of all possible
values of $f$ must be equal to unity, we have
\begin{equation}
  \int |a_f|^2\,df=1
  \tag{5.2}
\end{equation}
(analogously to relation (3.3) for the discrete spectrum).

Proceeding exactly as we did in deriving formula (3.5), and using the same
arguments, we write, on one hand,
\[
  \int\Psi\Psi^*\,dq=\int|a_f|^2\,df,
\]
and, on other hand,
\[
  \int\Psi\Psi^*\,dq
  =\iint a_f^*\Psi_f^*\Psi\,df\,dq.
\]
By comparing the two expressions we find the formula that determines the
expansion coefficients,
\begin{equation}
  a_f=\int\Psi(q)\Psi_f^*(q)\,dq,
  \tag{5.3}
\end{equation}
exactly analogous to (3.5).

To derive the normalization condition we now substitute (5.1) into (5.3):
\[
  a_f=\int a_{f'}\left(\int\Psi_{f'}\Psi_f^*\,dq\right)df'.
\]
This relation must hold for arbitrary $a_f$ and must therefore be satisfied
identically. For this to be so it is first of all necessary that the
coefficient of $a_{f'}$ under the integral sign (i.e. the integral
$\int\Psi_{f'}\Psi_f^*\,dq$) vanish for all $f'\ne f$. For $f'=f$ this
coefficient must become infinite (otherwise the integral over $df'$ would
simply be equal to zero). Thus the integral $\int\Psi_{f'}\Psi_f^*\,dq$ is a
function of the difference $f-f'$ which vanishes for non-zero values of the
argument and becomes infinite when the argument is zero. We denote this function
by $\delta(f'-f)$:
\begin{equation}
  \int\Psi_{f'}\Psi_f^*\,dq=\delta(f'-f).
  \tag{5.4}
\end{equation}

The manner in which the function $\delta(f'-f)$ becomes infinite at
$f'-f=0$ is determined by the requirement that there should be
\[
  \int\delta(f'-f)a_{f'}\,df'=a_f.
\]
\SourcePage{31}{34}
Clearly, for this to hold it is necessary that
\[
  \int\delta(f'-f)\,df'=1.
\]
A function defined in this way is called the
\textit{delta-function}\footnote{The delta-function was introduced into
theoretical physics by \textit{Dirac} (\textit{P.~A.~M. Dirac}).}. Let us write
out once more the formulas that define it. We have
\begin{equation}
  \delta(x)=0\quad\text{at }x\ne0,\qquad \delta(0)=\infty,
  \tag{5.5}
\end{equation}
and moreover in such a way that
\begin{equation}
  \int_{-\infty}^{+\infty}\delta(x)\,dx=1.
  \tag{5.6}
\end{equation}
As the limits of integration one may write any other pair between which the
point $x=0$ lies. If $f(x)$ is a function continuous at $x=0$, then
\begin{equation}
  \int_{-\infty}^{+\infty}\delta(x)f(x)\,dx=f(0).
  \tag{5.7}
\end{equation}
In more general form this formula can be written as
\begin{equation}
  \int\delta(x-a)f(x)\,dx=f(a),
  \tag{5.8}
\end{equation}
where the region of integration includes the point $x=a$, and $f(x)$ is
continuous at $x=a$. It is also obvious that the delta-function is even, i.e.
\begin{equation}
  \delta(-x)=\delta(x).
  \tag{5.9}
\end{equation}
Finally, writing
\[
  \int_{-\infty}^{\infty}\delta(\alpha x)\,dx
  =\int_{-\infty}^{\infty}\delta(y)\frac{dy}{|\alpha|}
  =\frac{1}{|\alpha|},
\]
we arrive at the conclusion that
\begin{equation}
  \delta(\alpha x)=\frac{1}{|\alpha|}\delta(x),
  \tag{5.10}
\end{equation}
where $\alpha$ is any constant.
\SourcePage{32}{35}
Formula (5.4) expresses the normalization rule for the eigenfunctions of the
continuous spectrum; it replaces condition (3.6) of the discrete spectrum. We
see that the functions $\Psi_f$ and $\Psi_{f'}$ with $f\ne f'$ remain
orthogonal to each other. The integrals of the squares $|\Psi_f|^2$ of the
functions of the continuous spectrum, however, diverge.

The functions $\Psi_f(q)$ satisfy one further relation, similar to (5.4). To
derive it we substitute (5.3) into (5.1), which gives
\[
  \Psi(q)=\int\Psi(q')
  \left(\int\Psi_f^*(q')\Psi_f(q)\,df\right)dq',
\]
whence we immediately conclude that one must have
\begin{equation}
  \int\Psi_f^*(q')\Psi_f(q)\,df=\delta(q-q').
  \tag{5.11}
\end{equation}
An analogous relation can, of course, be derived for the discrete spectrum,
where it takes the form
\begin{equation}
  \sum_n\Psi_n^*(q')\Psi_n(q)=\delta(q'-q).
  \tag{5.12}
\end{equation}

On comparing the pair of formulas (5.1) and (5.4) with the pair (5.3) and
(5.11), we see that, on the one hand, the functions $\Psi_f(q)$ effect the
expansion of the function $\Psi(q)$ with the expansion coefficients $a_f$, while,
on the other hand, formula (5.3) can be regarded as a completely analogous
expansion of the function $a_f\equiv a(f)$ in the functions $\Psi_f^*(q)$, with
$\Psi(q)$ playing the role of the expansion coefficients. The function $a(f)$,
like $\Psi(q)$, determines the state of the system completely; it is spoken of
as the wave function in the \textit{$f$-representation} (and the function
$\Psi(q)$ as the wave function in the \textit{$q$-representation}). Just as
$|\Psi(q)|^2$ determines the probability for the system to have coordinates in
a given interval $dq$, so $|a(f)|^2$ determines the probability of values of
the quantity $f$ in a given interval $df$. The functions $\Psi_f(q)$, for their
part, are, on the one hand, the eigenfunctions of the quantity $f$ in the
$q$-representation, while their complex conjugates $\Psi_f^*(q)$ serve as
eigenfunctions of the coordinate $q$ in the $f$-representation.

Suppose $\varphi(f)$ is a function of the quantity $f$, with $\varphi$ and
$f$ standing in one-to-one correspondence. Every function $\Psi_f(q)$ may then
equally be viewed as an eigenfunction of the quantity
$\varphi$. In doing so, however, the normalization of these functions must be
changed. Indeed, the eigen-
\SourcePage{33}{36}
functions $\Psi_\varphi(q)$ of the quantity $\varphi$ must be normalized by the
condition
\[
  \int\Psi_{\varphi(f')}\Psi_{\varphi(f)}^*\,dq
  =\delta[\varphi(f')-\varphi(f)],
\]
whereas the functions $\Psi_f$ are normalized by condition (5.4). The argument
of the delta-function vanishes at $f'=f$. For $f'$ close to $f$ we have
\[
  \varphi(f')-\varphi(f)=\frac{d\varphi(f)}{df}(f'-f).
\]
Making use of expression (5.10), one can write\footnote{In general, if
$\varphi(x)$ is a single-valued function (its inverse, however, may be
many-valued), then the formula
\begin{equation}
  \delta[\varphi(x)]
  =\sum_i\frac{1}{|\varphi'(\alpha_i)|}\delta(x-\alpha_i),
  \tag{5.13a}
\end{equation}
holds, where $\alpha_i$ are the roots of the equation $\varphi(x)=0$.}
\begin{equation}
  \delta[\varphi(f')-\varphi(f)]
  =\frac{1}{|d\varphi(f)/df|}\delta(f'-f).
  \tag{5.13}
\end{equation}
Comparison of (5.13) with (5.4) now shows that the functions $\Psi_\varphi$ and
$\Psi_f$ are related to each other by the relation
\begin{equation}
  \Psi_{\varphi(f)}
  =\frac{1}{\sqrt{|d\varphi(f)/df|}}\Psi_f.
  \tag{5.14}
\end{equation}

There exist physical quantities which, over some domain of their values,
possess a discrete spectrum, and over another possess a continuous one. For the
eigenfunctions of such a quantity, of course, all the same relations that were
derived in this and the preceding sections continue to hold. It need only be
noted that the complete system of functions is formed by the collection of the
eigenfunctions of both spectra together. Therefore the expansion of an
arbitrary wave function in the eigenfunctions of such a quantity has the form
\begin{equation}
  \Psi(q)=\sum_n a_n\Psi_n(q)+\int a_f\Psi_f(q)\,df,
  \tag{5.15}
\end{equation}
where the summation runs over the discrete spectrum, and the integration over
the entire continuous spectrum.

An example of a quantity possessing a continuous spectrum is the coordinate
$q$ itself. One readily sees that the associated operator reduces to
multiplication by $q$. Indeed,
\SourcePage{34}{37}
since the square $|\Psi(q)|^2$ governs the probability of different coordinate
values, the mean value of the coordinate is
\[
  \bar q=\int q|\Psi|^2\,dq=\int\Psi^*q\Psi\,dq.
\]
Comparing this expression with the definition of operators according to (3.8),
we see that\footnote{In what follows we shall agree, for simplicity of notation,
to write operators that reduce to multiplication by some quantity everywhere
simply as that quantity itself.}
\begin{equation}
  \hat q=q.
  \tag{5.16}
\end{equation}
The eigenfunctions of this operator must be determined, according to the general
rule, by the equation $q\Psi_{q_0}=q_0\Psi_{q_0}$, where $q_0$ is used
temporarily to denote concrete values of the coordinate, in distinction from
the variable $q$. Since this equality can be satisfied either when
$\Psi_{q_0}=0$ or when $q=q_0$, it is clear that the normalized eigenfunctions
are\footnote{The expansion coefficients of an arbitrary function $\Psi$ in these
eigenfunctions are $a_{q_0}=\int\Psi(q)\delta(q-q_0)\,dq=\Psi(q_0)$. The
probability of values of the coordinate in a given interval $dq_0$ is
$|a_{q_0}|^2\,dq_0=|\Psi(q_0)|^2\,dq_0$, as it should be.}
\begin{equation}
  \Psi_{q_0}=\delta(q-q_0).
  \tag{5.17}
\end{equation}
